\documentclass[10pt]{article}
\usepackage[margin=0.85in]{geometry}
\usepackage{amsmath}
\usepackage[hidelinks]{hyperref}

\begin{document}
\begin{center}
  {\Large\bfseries The thickness of complete graphs of prime order\par}
  \smallskip
  Proof of Conjecture 00000000163
\end{center}

Let $\theta(K_p)$ be the least number of planar layers covering the
edges of $K_p$. A cover gives a decomposition by assigning each edge to
one layer that contains it and deleting its other copies. Edge deletion
preserves planarity. We prove $\theta(K_p)=\lceil(p+2)/6\rceil$ for prime
$p$.

\section*{The edge-count lower bound}

The graphs $K_2$ and $K_3$ are planar and have edges, so each has
thickness one, as required. Now let $p\ge5$ be prime. Then $p$ is
divisible by neither two nor three, so $p=6m+1$ or $p=6m+5$.

If $t$ planar layers cover $K_p$, we may regard each layer as spanning
by adding isolated vertices. The simple planar edge bound from Euler's
formula gives at most $3(p-2)$ edges per layer. Counting the edges of
$K_p$ therefore gives
\[
 p(p-1)\le 6(p-2)t.
\]
If $p=6m+1$, then $m\ge1$ and
\[
 p(p-1)-6(p-2)m=12m>0.
\]
Thus $t\ge m+1$. If $p=6m+5$, then
\[
 p(p-1)-6(p-2)(m+1)=2>0,
\]
so $t\ge m+2$.

\section*{Matching planar covers}

Beineke and Harary construct a cover of $K_{6r+2}$ by $r+1$ planar
graphs for $r\ge1$~\cite[p.~855 and Figure~3]{bh}. Restricting a
planar cover of a complete graph to any subset of its vertices gives a
planar cover of the smaller complete graph.

For $p=6m+1$, restrict the cover of $K_{6m+2}$ to $p$ vertices. It has
$m+1$ layers. For $p=6m+5$, take $r=m+1$ and restrict the cover of
$K_{6m+8}$ to $p$ vertices. It has $m+2$ layers. These match the lower
bounds. Finally,
\[
 \left\lceil\frac{6m+3}{6}\right\rceil=m+1,\qquad
 \left\lceil\frac{6m+7}{6}\right\rceil=m+2,
\]
which gives the stated formula.

\section*{Lean verification}

The Lean 4 project checks the complete simple graph's edge count and the
list covering inequality, then proves the prime residue and ceiling
calculations. The planar Euler bound, the Beineke--Harary construction,
and preservation of planarity under restriction supply the stated
geometric inputs. Run \texttt{lake build} in the \texttt{lean} directory
to verify the deduction.

\begin{thebibliography}{1}
\bibitem{bh}
L.~W. Beineke and F.~Harary, \emph{The thickness of the complete graph},
Canadian Journal of Mathematics \textbf{17} (1965), 850--859.
\href{https://doi.org/10.4153/CJM-1965-084-2}{DOI 10.4153/CJM-1965-084-2}.
\end{thebibliography}

\end{document}
